% SPDX-License-Identifier: Apache-2.0
\section{Units, Processes and Parallelism}
\label{sec:e-units}

A unit is a struct that implements \code{crate::comp::Module}. The trait
has one method.

\begin{lstlisting}[caption={The unit trait. Inputs and outputs are type parameters, so one struct may implement it more than once.},label={lst:module}]
#[allow(async_fn_in_trait)]
pub trait Module<In, Out> {
    async fn run(&mut self, inputs: In, outputs: Out);
}
\end{lstlisting}

Inputs and outputs are type parameters rather than associated types, and
that choice buys something concrete: one struct may implement the trait
more than once, with different shapes, so a unit that works on 32-bit and
64-bit streams is one implementation and not two.

\subsection{Parallelism is stated, not implied}

Synthesis calls \code{run}. A unit with more than one concurrent process
starts one async function per process and joins them.

\begin{lstlisting}[caption={Parallelism written down. A unit with one process writes no join.},label={lst:join}]
impl Module<Ins, Outs> for DualPort {
    async fn run(&mut self, _i: Ins, _o: Outs) {
        join2(self.port_a(), self.port_b()).await;
    }
}
\end{lstlisting}

This replaces a rule with a construct. TxHDL said that several sequence
blocks in one module run concurrently, which was a property of having
declared them. It also contradicted itself: the specification stated that a
module holds one sequence, while three of its own modules declared
two~\cite{unification}. Here there is no rule to contradict. Parallelism is
a \code{join}, written where it happens, and a unit with one process does
not write one.

\subsection{Processes cannot take a mutable borrow}

The obvious way to write the two processes does not compile, and finding
that out cost ten lines.

\begin{lstlisting}[caption={What does not compile, and the error it gives. This is the code a Rust programmer writes first.},label={lst:e0499}]
async fn port_a(&mut self) { /* ... */ }
async fn port_b(&mut self) { /* ... */ }

// error[E0499]: cannot borrow `*self` as mutable
//               more than once at a time
join2(self.port_a(), self.port_b()).await;
\end{lstlisting}

Hardware processes share state: they drive registers. Rust permits one
mutable borrow at a time. The two facts collide directly.

The repair is that a register is a cell, and a process takes a shared
reference.

\begin{lstlisting}[caption={The repair. A register is a cell several processes reach, which is what a register is.},label={lst:reg}]
pub struct Reg<T: Copy> { v: core::cell::Cell<T> }

impl<T: Copy> Reg<T> {
    pub fn get(&self) -> T { self.v.get() }
    pub fn set(&self, x: T) { self.v.set(x) }
}

async fn port_a(&self) { self.hits_a.set(self.hits_a.get() + 1); }
\end{lstlisting}

\code{run} still takes \code{\&mut self}, and the processes it starts take
\code{\&self}. That compiles, and it also describes the hardware better: a
register is a thing several processes reach, which is what a register is.

\begin{pitfall}
This is the constraint most likely to surprise somebody writing a unit for
the first time, because the code that fails is the code a Rust programmer
would write first. It should be in the library's documentation and in the
error message, not only in a paper.
\end{pitfall}
